\chapter{Literature Review}\label{chap:background}

This chapter surveys the academic literature and technical background that underpins the
design decisions made in Task~Flow.
The review spans seven areas: existing task-management tools and their limitations,
desktop application frameworks, offline-first architecture patterns, hybrid database
designs, real-time communication protocols, authentication and security, and the
evolution of Large Language Models (LLMs) as productivity aids.
It concludes with a brief treatment of project-visualisation paradigms.

% ─────────────────────────────────────────────────────────────────────────────
\section{Task Management and Productivity Software}\label{sec:bg_tasktools}
% ─────────────────────────────────────────────────────────────────────────────

Task management software encompasses a broad spectrum of tools---from simple personal
to-do lists to enterprise project-management suites.
The most widely cited personal productivity framework is David Allen's
\emph{Getting Things Done}~(GTD)~\cite{allen2001gtd}, which argues that externalising
tasks into a trusted system frees cognitive resources for execution rather than recall.
Modern software tools attempt to operationalise this principle in varying ways.

\textbf{Trello} implements the Kanban paradigm~\cite{anderson2010kanban} using cards
on boards with user-defined columns.
Its simplicity has made it popular for personal and small-team use, but it provides
no offline capability; every drag-and-drop operation requires an active connection.
\textbf{Asana} and \textbf{Jira} extend the model with structured project hierarchies,
sprint planning, and reporting dashboards, but both are unambiguously cloud-first and
offer only read-only offline caching at best.
\textbf{Notion} introduces a flexible block-based editor that unifies notes, databases,
and tasks in one surface, yet it shares the same connectivity dependency.

Table~\ref{tab:tool_comparison} summarises the key feature dimensions of these
tools relative to Task~Flow.

\begin{table}[H]
\centering
\caption{Feature comparison of popular task-management tools vs.\ Task~Flow.}
\label{tab:tool_comparison}
\begin{tabular}{lcccccc}
\toprule
\textbf{Feature} & \textbf{Trello} & \textbf{Asana} & \textbf{Jira} & \textbf{Notion} & \textbf{Task~Flow} \\
\midrule
Offline operation      & \texttimes & \texttimes & \texttimes & \texttimes & \checkmark \\
Local persistence      & \texttimes & \texttimes & \texttimes & Limited    & \checkmark \\
Transparent sync       & \texttimes & \texttimes & \texttimes & \texttimes & \checkmark \\
Embedded AI assistant  & Limited    & Limited    & Limited    & Limited    & \checkmark \\
Desktop installer      & \texttimes & \texttimes & \texttimes & \checkmark & \checkmark \\
Kanban view            & \checkmark & \checkmark & \checkmark & \checkmark & \checkmark \\
Gantt / timeline view  & \texttimes & \checkmark & \checkmark & \texttimes & \checkmark \\
Calendar view          & \texttimes & \checkmark & \checkmark & \checkmark & \checkmark \\
Integrated messaging   & \texttimes & \texttimes & \texttimes & \texttimes & \checkmark \\
\bottomrule
\end{tabular}
\end{table}

The pattern that emerges from this survey is consistent: each tool excels in one or two
dimensions but requires multiple third-party integrations to achieve a comparable feature
breadth to Task~Flow, and none treats offline reliability as a design priority.

% ─────────────────────────────────────────────────────────────────────────────
\section{Desktop Application Frameworks}\label{sec:bg_desktop}
% ─────────────────────────────────────────────────────────────────────────────

The choice of desktop application framework has long-term implications for developer
productivity, binary size, platform reach, and the ability to reuse existing web
technology investments.

\textbf{Native frameworks} such as WinForms, WPF (Windows), Cocoa (macOS), and GTK
(Linux) offer tight OS integration and the best runtime performance but impose a
per-platform development effort and cannot reuse the React/TypeScript frontend code
already written for Task~Flow.

\textbf{Tauri}~\cite{electrondocs2024} is a newer alternative to Electron that renders
web content in the system's native WebView and uses Rust for the application shell,
resulting in significantly smaller binaries.
However, its Rust back-end integration would require a full rewrite of the .NET server
component, and its ecosystem maturity is substantially lower than Electron's.

\textbf{Electron}~\cite{electrondocs2024}, developed by GitHub and now maintained by
the OpenJS Foundation, embeds a Chromium browser and a Node.js runtime in a single
cross-platform binary.
The web content layer communicates with the Node.js main process via inter-process
communication (IPC).
This architecture allows the complete React~+~TypeScript frontend to be reused without
modification.
For Task~Flow, the .NET backend is launched as a child process by the Electron main
process, with port negotiation performed via a token on standard output
(\texttt{TASKFLOW\_BACKEND\_READY:\{address\}}).
The resulting installer sizes are larger than Tauri equivalents but remain well within
acceptable bounds for a desktop application.
Electron's comprehensive documentation, wide community adoption, and first-class support
in \texttt{electron-builder}~\cite{electrondocs2024} made it the optimal choice for
this project.

% ─────────────────────────────────────────────────────────────────────────────
\section{Offline-First Architecture}\label{sec:bg_offline}
% ─────────────────────────────────────────────────────────────────────────────

The term \emph{offline-first} was popularised in the context of progressive web
applications and mobile clients, but the underlying principle applies equally to
desktop software: \emph{design the system to function fully without network access,
and treat synchronisation as an enhancement rather than a prerequisite}.

Kleppmann's \emph{Designing Data-Intensive Applications}~\cite{kleppmann2017designing}
provides a thorough treatment of the trade-offs involved.
In the context of the CAP theorem~\cite{brewer2012cap}, an offline-first system
prioritises \emph{Availability} and \emph{Partition Tolerance} over strict
\emph{Consistency}: local writes are always accepted, and consistency is restored
asynchronously when connectivity returns.

\subsection*{The Transactional Outbox Pattern}

A key challenge in offline-first systems is ensuring that locally committed changes are
eventually propagated to the remote store without duplication or loss.
The \emph{transactional outbox pattern}~\cite{fowler2002patterns} addresses this by
writing both the business entity change and a corresponding outbox entry in the same
local database transaction.
A background worker polls the outbox and replays pending entries to the remote system,
advancing each entry through a state machine: \texttt{Pending} $\to$
\texttt{Processing} $\to$ \texttt{Synced} (or \texttt{Failed} after exhausting retries).
This design guarantees at-least-once delivery while keeping the application logic
decoupled from network availability.

\subsection*{Conflict Resolution}

When multiple clients modify the same entity while offline, the system must resolve
divergent versions upon reconnection.
Conflict-free Replicated Data Types (CRDTs)~\cite{shapiro2011crdt,kleppmann2019crdt}
provide a mathematically principled approach: data structures are designed so that
any two states can be merged without coordination.
Implementing CRDTs in a production system adds significant complexity; the current
Task~Flow implementation adopts a simpler \emph{last-write-wins} strategy (using the
\texttt{UpdatedAt} timestamp as the arbitration key), with CRDTs noted as a long-term
enhancement.

% ─────────────────────────────────────────────────────────────────────────────
\section{Hybrid Database Architectures}\label{sec:bg_databases}
% ─────────────────────────────────────────────────────────────────────────────

\subsection*{SQLite as a Local Store}

SQLite~\cite{hipp2022sqlite} is a self-contained, serverless, zero-configuration,
transactional SQL database engine embedded as a library in the host application.
It writes to a single file on disk, requires no background daemon process, and
supports full ACID transactions.
These properties make it the natural choice for a local persistence layer: the database
travels with the application, survives process restarts, and performs all reads and
writes without network round-trips.
At the time of writing, SQLite is the most widely deployed database engine in the world,
present in mobile operating systems, browsers, and billions of embedded devices.

In Task~Flow, SQLite is the \emph{single source of truth}.
All application features read from and write to SQLite first; MongoDB involvement is
entirely asynchronous and does not block the user-visible operation path.

\subsection*{MongoDB as a Cloud Relay}

MongoDB~\cite{chodorow2013mongodb} is a document-oriented database that stores data as
BSON documents grouped in collections.
Its flexible schema accommodates evolving data shapes without migration overhead and
its horizontal scaling capabilities make it suitable for cloud deployment.

In Task~Flow, MongoDB serves as a \emph{synchronisation relay} for the subset of data
that must be visible across devices or shared among team members: mirrored tasks,
notifications, team invitations, and user-presence indicators.
The document \texttt{\_id} for each mirrored document is the entity's \texttt{SyncId}
(a UUID generated at creation time), ensuring idempotent upserts and stable cross-device
identity.
This dual-database pattern---ACID-safe local SQLite as primary, schemaless MongoDB as
asynchronous relay---is a pragmatic alternative to full CRDT-based synchronisation when
conflict probability is low and eventual consistency is acceptable.

\subsection*{Entity Framework Core}

Entity Framework Core~\cite{efcoredocs2024} is the official Object-Relational Mapper
(ORM) for .NET.
It provides a LINQ-based query API, change tracking, and migration tooling that
translates strongly typed C\# entity classes into SQL DDL and DML statements.
Its provider model allows the same application code to target SQLite in development and
testing and PostgreSQL or SQL~Server in production deployment without code changes---a
property exploited via the \texttt{DatabaseProviderConfig} switch in Task~Flow's
\texttt{Startup.cs}.

% ─────────────────────────────────────────────────────────────────────────────
\section{Real-Time Communication Patterns}\label{sec:bg_realtime}
% ─────────────────────────────────────────────────────────────────────────────

Modern collaborative applications require the server to \emph{push} state changes to
connected clients rather than waiting for clients to poll.
Three primary mechanisms are relevant to Task~Flow.

\subsection*{WebSockets}

The WebSocket protocol~\cite{rfc6455} establishes a persistent, full-duplex TCP
connection between a client and server after an HTTP upgrade handshake.
Both parties may send frames at any time without the overhead of repeated HTTP headers.
WebSockets are well-suited for low-latency bidirectional communication such as chat
messages and live presence indicators.

\subsection*{SignalR}

ASP.NET Core SignalR~\cite{signalrdocs2024} is a higher-level abstraction over
WebSockets, Server-Sent Events, and long polling.
It provides a hub-based RPC model in which the server can invoke named methods on
connected clients, and clients can invoke methods on the server.
SignalR automatically negotiates the best available transport and handles reconnection.
Task~Flow uses SignalR's \texttt{NotificationHub} to broadcast real-time events:
connectivity-state changes, synchronisation progress, task reminders, and new
notification payloads are all delivered to the React client via a single hub connection.

\subsection*{Server-Sent Events}

Server-Sent Events (SSE) is a unidirectional streaming protocol built on standard
HTTP: the server holds the response connection open and writes \texttt{text/event-stream}
chunks as data becomes available~\cite{signalrdocs2024}.
SSE is the appropriate choice for AI response streaming because the communication is
inherently one-directional (server to client) and does not require the connection upgrade
ceremony of WebSockets.
Task~Flow's \texttt{ChatbotController} exposes a \texttt{POST
/api/chatbot/conversations/\{id\}/messages/stream} endpoint that writes Mistral API
response chunks directly to the SSE response as they arrive, enabling the React
frontend to render the AI reply token by token.

% ─────────────────────────────────────────────────────────────────────────────
\section{Authentication and Security}\label{sec:bg_security}
% ─────────────────────────────────────────────────────────────────────────────

\subsection*{JSON Web Tokens}

JSON Web Tokens (JWT)~\cite{rfc7519} are a compact, URL-safe means of representing
claims between two parties.
A JWT consists of three Base64URL-encoded sections---header, payload, and
signature---concatenated with dots.
The signature, computed using HMAC-SHA256 (HS256) over the header and payload,
allows the receiving party to verify the token's integrity without querying a
central session store.
This stateless verification model is particularly valuable in a desktop application
context where the backend is a local process: there is no shared session infrastructure
between restarts.
Task~Flow issues JWTs with a 24-hour expiry; the signing key is injected at runtime
via the \texttt{TASKFLOW\_JWT\_KEY} environment variable, falling back to a random
32-byte hex key generated at startup if none is provided.

\subsection*{Password Hashing with bcrypt}

Storing passwords as plaintext or with reversible encryption is a critical security
failure.
The bcrypt algorithm~\cite{provos1999bcrypt} applies an adaptive key-derivation function
with a configurable cost factor, deliberately making hash computation slow.
As hardware performance improves, the cost factor can be increased to maintain the same
brute-force resistance.
Task~Flow uses BCrypt.Net-Next with the default cost factor to hash passwords at
registration and verify them at login, ensuring that even a full database compromise
does not directly expose user credentials.

\subsection*{OWASP Top 10 Considerations}

The Open Web Application Security Project (OWASP) publishes an annually updated list of
the ten most critical web application security risks~\cite{owasp2021}.
The most relevant to Task~Flow are:
\emph{Broken Access Control}---mitigated by validating the JWT \texttt{sub} claim on
every protected endpoint and checking entity ownership before returning or modifying
data;
\emph{Cryptographic Failures}---addressed by bcrypt hashing and HS256 JWT signing;
\emph{Injection}---mitigated by EF Core's parameterised queries which eliminate raw SQL
string concatenation; and
\emph{Security Misconfiguration}---acknowledged in the context of the Mistral API key,
which should be supplied via an environment variable rather than embedded in source code.

\subsection*{Role-Based Access Control}

Ferraiolo and Kuhn's foundational work on role-based access
control~(RBAC)~\cite{sandhu1996rbac} defines a model in which permissions are assigned
to roles and roles are assigned to users, rather than assigning permissions directly to
users.
Task~Flow implements a lightweight RBAC model for team membership: each
\texttt{TeamMember} record carries a \texttt{Role} field, and the Teams controller
enforces that only the team owner may perform administrative operations such as removing
members or disbanding the team.

% ─────────────────────────────────────────────────────────────────────────────
\section{Large Language Models and AI-Assisted Productivity}\label{sec:bg_ai}
% ─────────────────────────────────────────────────────────────────────────────

\subsection*{The Transformer Architecture}

The transformer architecture, introduced by Vaswani et~al.~\cite{vaswani2017attention}
in the seminal paper \textit{Attention Is All You Need}, replaced recurrent
sequence-to-sequence models with a purely attention-based mechanism.
The self-attention operation allows each token in a sequence to attend to every other
token in parallel, enabling efficient training on large corpora and strong performance
on a wide range of language tasks.
This architecture became the foundation for all major LLMs in subsequent years.

\subsection*{GPT-3 and the Emergence of Instruction Following}

Brown et~al.'s GPT-3~\cite{brown2020gpt3} demonstrated that scaling transformer models
to 175 billion parameters produced qualitatively new capabilities, including few-shot
task completion from natural-language prompts.
The subsequent development of instruction-tuned variants established the prompt-response
interaction model that underpins modern conversational AI assistants.

\subsection*{Mistral Models}

Jiang et~al.~\cite{jiang2023mistral} introduced Mistral~7B, a 7-billion-parameter
open-weight model that outperformed LLaMA~2~13B on several benchmarks despite being
half the size, achieved through grouped-query attention and sliding-window attention
mechanisms.
Task~Flow integrates the Mistral API at three model tiers:
\texttt{mistral-small-latest} for general conversational use;
\texttt{codestral-latest} for the code-generation mode, a model fine-tuned on a large
corpus of programming languages; and
\texttt{mistral-ocr-latest} for the document-scanner mode, which accepts base64-encoded
images and PDF files and returns extracted and structured text.
The 20-message conversation history window provides conversational context while
bounding the per-request token count.

\subsection*{AI-Assisted Productivity Research}

Recent studies examining AI coding assistants report productivity gains of 26--55\%
on specific programming tasks, with larger gains on routine boilerplate generation
and smaller gains on complex architectural reasoning.
The integration of an AI assistant directly within the task-management interface---rather
than requiring the user to switch to a separate tool---reduces context-switching overhead
and allows the assistant to be consulted in the flow of work.

% ─────────────────────────────────────────────────────────────────────────────
\section{Project Visualisation Paradigms}\label{sec:bg_viz}
% ─────────────────────────────────────────────────────────────────────────────

Different users hold different mental models of their work, and no single visualisation
paradigm suits all cognitive styles or project types.
A mature task-management tool should offer multiple views over the same underlying data.

\textbf{Kanban boards} originated in Toyota's production system as physical cards
(\emph{kanban} = ``signal card'' in Japanese) that controlled the flow of work between
manufacturing stages~\cite{anderson2010kanban}.
Their digital equivalent presents tasks as cards in columns representing workflow stages,
giving teams an immediate visual signal of bottlenecks (columns with too many cards)
and idle capacity.
Drag-and-drop interaction directly maps the physical act of moving a card from one column
to the next.

\textbf{Gantt charts} were developed by Henry Gantt in the early twentieth century as
project-schedule representations that map tasks onto a horizontal time axis~\cite{wilson2003gantt}.
They are particularly effective for deadline-sensitive planning because the chart
immediately communicates which tasks are in progress, which are complete, and which are
overdue relative to a calendar.

\textbf{Calendar views} represent tasks as events on a time grid, a mental model
familiar to users from personal calendars.
A weekly or monthly grid makes it easy to reason about task density and identify
overloaded days.

\textbf{Tabular views} surface all task attributes as a structured grid, enabling rapid
scanning, sorting, and filtering without the spatial metaphor of Kanban columns or
timeline bars.
This view is preferred by users who think in terms of data rather than process flow.

Task~Flow provides all four of these paradigms plus a smart-grouped default list
(grouping tasks by temporal proximity: Overdue, Today, This Week, Later, Completed),
giving each user the freedom to work in the style that best fits their workflow at any
given moment.

% ─────────────────────────────────────────────────────────────────────────────
\section{Summary}\label{sec:bg_summary}
% ─────────────────────────────────────────────────────────────────────────────

The literature review reveals a clear pattern of unmet needs in the current landscape of
task-management software.
Existing tools are built on cloud-first assumptions that render them fragile in
low-connectivity environments.
The offline-first literature, particularly the transactional outbox pattern described
by Fowler~\cite{fowler2002patterns} and the data-intensive application principles
formalised by Kleppmann~\cite{kleppmann2017designing}, provides the theoretical basis
for a reliable local-first persistence and synchronisation subsystem.
The dual-database architecture combining SQLite and MongoDB is a pragmatic application
of these principles that avoids the complexity of full CRDT-based conflict resolution
while still enabling cross-device data relay.

In the domain of AI integration, the Mistral family of models~\cite{jiang2023mistral}
offers an open-weight, commercially accessible set of capabilities---general language,
code generation, and OCR---that align precisely with the three modes required by the
system.
Server-Sent Events provide the appropriate transport mechanism for streaming AI
responses to the client with minimal protocol overhead.

Electron's maturity and React-reuse properties make it the most pragmatic desktop
packaging choice, while JWT and bcrypt together provide a well-understood and
academically grounded authentication and security foundation.

The following chapter translates these theoretical foundations into a concrete system
design, detailing the architecture, data model, and implementation of each subsystem
in Task~Flow.
